%! TEX root = main.tex

\begin{exercise}\label{ex:em-1}
Charges $q_1 = +2.0\ \mu\mathrm{C}$ at $x = 0$ and $q_2 = -3.0\ \mu\mathrm{C}$ at $x = 40\ \mathrm{cm}$ are fixed on the $x$-axis. A third charge $q_3 = +5.0\ \mu\mathrm{C}$ is placed at $x = 15\ \mathrm{cm}$.
\begin{enumerate}[label=(\alph*)]
  \item Find the force on $q_3$ due to $q_1$ and the force on $q_3$ due to $q_2$: magnitude and direction of each.
  \item Find the net force on $q_3$ (magnitude and direction).
  \item At what $x$ between the two fixed charges does the electric field vanish?
\end{enumerate}
\end{exercise}

\begin{exercise}\label{ex:em-2}
Two large parallel plates are $3.0\ \mathrm{mm}$ apart and connected to a $400\ \mathrm{V}$ supply, so the field between them is uniform. An electron is released from rest at the negative plate.
\begin{enumerate}[label=(\alph*)]
  \item Find the electric field strength between the plates.
  \item Find the force on the electron and its acceleration.
  \item Find the electron's speed when it reaches the positive plate, using energy.
\end{enumerate}
\end{exercise}

\begin{exercise}\label{ex:em-3}
A parallel-plate capacitor has plate area $50\ \mathrm{cm^2}$, separation $0.50\ \mathrm{mm}$, and a dielectric with $\varepsilon_r = 4.0$ between the plates.
\begin{enumerate}[label=(\alph*)]
  \item Find its capacitance.
  \item Find the stored charge and energy when it is connected to $24\ \mathrm{V}$.
  \item The battery is disconnected and the plates are then pulled to twice their original separation. State what happens to $C$, $U$, and the stored energy, and explain where the extra work goes.
\end{enumerate}
\end{exercise}

\begin{exercise}\label{ex:em-4}
A $12\ \mathrm{V}$ battery drives $R_1 = 150\ \Omega$ in series with the parallel combination of $R_2 = 220\ \Omega$ and $R_3 = 330\ \Omega$.
\begin{enumerate}[label=(\alph*)]
  \item Find the total resistance seen by the battery and the current it delivers.
  \item Find the current through each resistor and the voltage across each.
  \item Find the power dissipated in $R_1$ and the total power delivered by the battery.
\end{enumerate}
\end{exercise}

\begin{exercise}\label{ex:em-5}
A proton enters a uniform magnetic field of $B = 0.25\ \mathrm{T}$ at $v = 4.0\times 10^{5}\ \mathrm{m/s}$, perpendicular to the field.
\begin{enumerate}[label=(\alph*)]
  \item Determine the direction of the initial force (describe it in words, choosing axes yourself).
  \item Find the radius and the period of the resulting circular path.
  \item An electron enters with the same velocity in the same field. How do its radius and period compare with the proton's? Justify from the formulas.
\end{enumerate}
\end{exercise}

\begin{exercise}\label{ex:em-6}
Prove that a magnetic field does no work on a moving charge. Start from $\vec{F} = q\vec{v}\times\vec{B}$ and the definition of mechanical work, and deduce that a particle moving in a pure magnetic field keeps constant speed and constant kinetic energy.
\end{exercise}

\begin{exercise}\label{ex:em-7}
A metal rod of length $L = 50\ \mathrm{cm}$ slides without friction along horizontal rails in a uniform vertical field $B = 0.20\ \mathrm{T}$, at constant speed $v = 4.0\ \mathrm{m/s}$. The rails are connected to a resistor $R = 5.0\ \Omega$; the rod and rails have negligible resistance.
\begin{enumerate}[label=(\alph*)]
  \item Starting from the magnetic force on the charge carriers in the rod, derive that the induced voltage across the moving rod is $U_\text{ind} = BLv$, and state the condition under which this formula holds.
  \item Find the induced voltage, the induced current, and its direction (explain your reasoning with Lenz's law).
  \item Find the external force needed to keep the rod at constant speed, and the mechanical power supplied. Compare with the power dissipated in the resistor.
\end{enumerate}
\end{exercise}

\begin{exercise}\label{ex:em-8}
A proton moves in a uniform magnetic field $\vec{B} = (0,0,B)$ in vacuum, with initial velocity in the $xy$-plane. Its motion obeys
\[
  \frac{\mathrm{d}\vec{v}}{\mathrm{d}t} = \frac{q}{m}\,\vec{v}\times\vec{B}.
\]
Simulate this numerically in a language of your choice.
\begin{enumerate}[label=(\alph*)]
  \item Implement the semi-implicit Euler step $\vec{v}_{n+1} = \vec{v}_n + \frac{q}{m}\,\Delta t\,(\vec{v}_n\times\vec{B})$, then $\vec{r}_{n+1} = \vec{r}_n + \Delta t\,\vec{v}_{n+1}$, and also plain (explicit) Euler where $\vec{r}$ is updated with $\vec{v}_n$.
  \item Run both for a few periods with $\Delta t = T/100$ (estimate $T$ from \cref{exm:proton-circle}) and plot the orbits. Measure how the particle's speed drifts over one period for each scheme.
  \item Compare the simulated orbit radius with the theoretical $r = mv/(qB)$ and explain any discrepancy. State which scheme you would use in a real physics engine and why.
\end{enumerate}
\end{exercise}
